\documentclass[UTF8,a4paper,12pt]{ctexart}
\usepackage[margin=2.5cm]{geometry}
\usepackage{setspace}
\usepackage{enumitem}
\usepackage{tikz}
\usetikzlibrary{arrows.meta,positioning,calc}
\usepackage{hyperref}
\hypersetup{hidelinks}
\setstretch{1.35}
\setlength{\parindent}{2em}
\setlength{\parskip}{0pt}
\ctexset{section={format=\large\bfseries,aftername=\quad},subsection={format=\normalsize\bfseries,aftername=\quad}}
\newcommand{\refmark}[1]{\textsuperscript{[#1]}}
\begin{document}
\title{室内高精度定位技术研究综述}
\author{}
\date{}
\maketitle
\vspace{-3em}
\begin{abstract}
随着物联网（IoT）和位置服务（LBS）在智能制造、仓储物流及日常生活中的深度普及，室内定位技术已成为学术界和工业界的研究热点。由于传统的全球导航卫星系统（GNSS）在室内环境中存在严重的信号遮挡与衰减，难以提供有效的空间位置信息\refmark{1}。本文系统梳理了室内定位技术的发展历程，重点剖析了从早期基于信号强度的粗放定位向基于到达角（AoA）、超宽带（UWB）等高精度定位演进过程中的核心技术与算法演变。特别地，本文探讨了面向物联网低成本需求的前沿技术，深入论证了基于分布式相控阵的“REV 相位反演法”结合“二维 MUSIC 算法”在实现低成本、高精度定位中的巨大潜力。最后，本文对多源传感器融合及通导一体化的未来发展趋势进行了展望。
\end{abstract}

\section{引言}
位置信息是连接物理世界与数字信息世界的关键桥梁。尽管 GPS 和北斗等卫星导航系统在室外开阔环境下能够提供高精度的定位服务，但人类日常活动有极高比例是在室内进行的。GNSS 信号由于无法穿透建筑物墙体，在室内环境中基本失效\refmark{2}。为了打破室内定位的瓶颈，研究人员探索了诸如 Wi-Fi、蓝牙、超宽带（UWB）、5G 蜂窝网络以及地磁和视觉等多种技术途径\refmark{3}。近年来，随着工业物联网和智能场景的复杂化，传统的米级定位精度已无法满足精准追踪、机器人导航等场景的需求，厘米级到亚米级的高精度室内定位技术由此迎来了迅猛发展并逐渐成为行业的主流演进方向\refmark{4}。

\section{传统室内定位技术及局限性}
早期的室内定位主要依赖于现有的通信基础设施，通过测量射频信号的物理特征来实现位置估算。
\subsection{基于 RSSI 的测距与指纹定位}
接收信号强度指示（RSSI）是早期 Wi-Fi 和蓝牙网络中最普遍的定位信息源。其基本原理是利用电磁波在空间中的对数距离衰减模型，将信号强度转换为物理距离。然而，室内环境存在严重的墙体遮挡和多径效应，导致 RSSI 信号产生剧烈波动。为了克服这一问题，位置指纹定位（Fingerprinting）应运而生。该方法分为离线建库和在线匹配两个阶段：在定位区域内密集采集 RSSI 建立特征指纹库，随后使用 K 近邻（KNN）或加权 K 近邻（WKNN）等算法进行匹配\refmark{6}。尽管引入了滑动平均、中值滤波等数据预处理手段来降低噪声，且指纹法无需建立复杂的信道模型，但其离线建库耗费大量人力，且极易受到环境动态变化（如人员走动、物体搬移）的干扰，定位精度通常只能维持在 3 至 5 米之间，难以突破高精度的瓶颈\refmark{6}。
\begin{figure}[htbp]
\centering
\begin{tikzpicture}[>=Latex, node distance=9mm and 11mm, every node/.style={font=\small}, box/.style={draw=black!60,rounded corners=2pt,align=center,minimum height=10mm,minimum width=24mm,fill=blue!4}]
\node[box] (grid) {参考点\\采集 RSSI};
\node[box,right=of grid] (db) {离线指纹库};
\node[box,below=13mm of grid] (live) {待定位设备\\测量 RSSI};
\node[box,right=of live] (match) {KNN/WKNN\\指纹匹配};
\node[box,right=of match] (pos) {位置估计};
\draw[->,thick] (grid)--(db);
\draw[->,thick] (live)--(match);
\draw[->,thick] (db)--(match);
\draw[->,thick] (match)--(pos);
\node[font=\footnotesize,above=2mm of grid] {离线阶段};
\node[font=\footnotesize,below=2mm of live] {在线阶段};
\end{tikzpicture}
\caption{RSSI 指纹定位的离线建库与在线匹配流程}
\label{fig:rssi}
\end{figure}
\subsection{传统低功耗蓝牙（BLE）定位}
在蓝牙 5.1 标准发布之前，基于 iBeacon 的低功耗蓝牙定位主要也是通过广播包的 RSSI 进行测距。虽然其功耗低、部署灵活，但与 Wi-Fi 类似，定位精度受限于多径衰落，无法满足严苛的工业级需求\refmark{1}。

\section{高精度室内定位的核心算法原理}
为实现亚米级乃至厘米级的高精度定位，底层算法模型开始从测量信号“强度”向测量信号的“时间”与“空间角度”转变。
\subsection{基于时间的高精度测距（TOA/TDOA）}
飞行时间（TOA）算法通过精确测量无线电信号单程飞行时间来解算距离，而到达时间差（TDOA）则利用多基站接收同一信号的时间差进行双曲线方程解算。TDOA 的优势在于待测标签无需与基站进行严苛的时间同步。在 UWB 等大带宽通信系统中，基于时间的算法能够实现极高的空间分辨率\refmark{5}。
\subsection{空间角度定位：到达角（AoA）与多天线阵列}
到达角（AoA）定位技术通过测量无线电信号入射到接收端多天线阵列不同阵元上的相位差，进而换算为信号的到达角度，最后结合空间几何关系进行交叉定位。AoA 方法不需要高精度的时钟同步，抗干扰能力极强，是当前实现物联网高精度定位的核心算法基础\refmark{8}。针对多径环境对 AoA 估计精度的影响，研究表明通过推导 AoA 估计误差的渐近方差，为不同精度的估计角度分配动态权重（Weighted AoA），可将系统的中位数定位误差显著降低 20\%\refmark{7}。
\begin{figure}[htbp]
\centering
\begin{tikzpicture}[>=Latex,scale=0.9,every node/.style={font=\small}]
\coordinate (A) at (-2.8,0);
\coordinate (B) at (2.8,0);
\coordinate (T) at (0.35,2.55);
\draw[dashed,gray] (-3.5,0)--(3.5,0);
\draw[very thick,blue!70!black,-Latex] (T)--(A);
\draw[very thick,blue!70!black,-Latex] (T)--(B);
\foreach \x in {-3.10,-2.80,-2.50,2.50,2.80,3.10}{\filldraw[fill=blue!35,draw=black] (\x,0) circle (0.075);}
\filldraw[fill=black] (T) circle (0.075);
\node[above=2mm of T] {待定位节点};
\node[below=5mm of A] {接收阵列 1};
\node[below=5mm of B] {接收阵列 2};
\draw[thick] ($(A)+(0.67,0)$) arc[start angle=0,end angle=39,radius=0.67];
\draw[thick] ($(B)+(-0.67,0)$) arc[start angle=180,end angle=134,radius=0.67];
\node at (-1.92,0.32) {$\theta_1$};
\node at (1.93,0.36) {$\theta_2$};
\end{tikzpicture}
\caption{双接收阵列到达角交会定位示意}
\label{fig:aoa}
\end{figure}

\section{面向物联网的高精度定位技术突破}
在算法理论的支撑下，主流通信硬件在底层协议和天线设计层面进行了深刻的革新。
\subsection{UWB 与 5G 使能的泛在定位}
超宽带（UWB）技术采用纳秒级的极窄脉冲信号进行通信，具有极强的穿透能力和多径分离能力。在视距（LOS）环境下，UWB 定位精度可达厘米级（通常为 10--30 厘米），是目前成熟度最高的高精度方案之一\refmark{5}。然而，UWB 专用硬件成本较高，限制了其在海量低成本物联网设备中的普及。与此同时，5G 蜂窝网络凭借其大带宽和大规模 MIMO 天线阵列，为室内外无缝定位提供了新思路。利用 5G 共频带信号（TC-OFDM）等技术，高精度定位正逐步演变为通信网络基础设施的内生能力，向“在线即在位”的目标演进\refmark{4}。
\subsection{蓝牙 AoA 技术与算法优化}
蓝牙 5.1 核心规范引入了寻向功能（Direction Finding），使得蓝牙定位跨入了高精度时代。利用采集到的连续音延展（CTE）数据包的 I/Q 样本，系统能够计算出信号到达角。针对接收机在单天线阵元间切换所产生的开关噪声及多径干扰，采用非线性递归最小二乘法（NRLS）与无迹卡尔曼滤波（UKF）相融合的处理架构，可使 AoA 估计误差平均降低 3.9 度，大幅提升系统的鲁棒性\refmark{8}。
\subsection{分布式相控阵使能的 REV+MUSIC 定位前沿}
在 AoA 角度估计算法中，经典的二维多重信号分类算法（2D MUSIC）具备出色的超分辨率特性，但其传统的全平面空间谱峰搜索带来了极高的计算复杂度，难以直接应用于算力受限的物联网边缘节点\refmark{9}。为解决算力瓶颈，基于多项式求根替代谱峰搜索的快速降维 2D MUSIC 算法被提出，大幅提升了解算效率\refmark{9}。

更为重要的是，针对物联网设备对“低成本、高精度”的核心诉求，基于分布式相控阵使能的联合技术路线展现出了卓越的潜力。该方案创新性地采用了“REV（旋转元电场矢量）相位反演法”来获取低成本的相位信息，并与优化的“二维 MUSIC 算法”深度结合以实现高分辨率的角度估计。通过构建从信道模拟到定位解算的全链路算法仿真，该 REV+MUSIC 联合方案在单点信源定位实验中，能够将平均定位误差稳定控制在 20cm 至 30cm 范围内。此外，结合 FEKO 等专业电磁仿真软件对系统物理布局进行贴近现实环境的迭代优化，进一步验证了该分布式相控阵系统在复杂室内环境下的鲁棒性与高精度性能，为海量物联网节点的低成本厘米级定位开辟了全新的技术路径。

\section{挑战与未来演进展望}
尽管高精度室内定位技术取得了长足进步，但在规模化商业落地中仍面临诸多挑战。未来的演进方向将主要聚焦于以下几个方面：
\begin{enumerate}[leftmargin=2em]
\item 非视距（NLOS）误差抑制：复杂的室内遮挡会导致信号直射径丢失，带来巨大的测距与测角误差。利用机器学习进行 NLOS 鉴别与误差动态补偿，将是提升系统稳定性的关键\refmark{1}。
\item 多源异构传感器融合：单一通信射频技术难以应对所有极端场景。未来的发展趋势是将无线射频信号（AoA/UWB）与物联网设备内置的惯性测量单元（IMU）、视觉传感器及地磁进行深度融合滤波（如扩展卡尔曼滤波），实现优势互补\refmark{5}。
\item 通导一体化网络融合：未来的室内定位将不再依赖大量外挂的独立基站，而是与 5G/6G 通信网络、智能照明（VLC 可见光通信）等基础设施深度结合，实现通信与高精度导航定位的一体化协同服务\refmark{2}\refmark{4}。
\end{enumerate}

\section{结论}
室内定位技术正经历从基于 RSSI 信号强度的粗放式估算，向以 AoA 空间测角和广义到达时间差为核心的高精度阶段的深刻变革。面对物联网时代对海量终端设备低成本与高精度的双重苛刻要求，诸如蓝牙多天线阵列优化以及“REV 相位反演+MUSIC”的分布式相控阵定位系统，正成为填补技术空白的关键力量。未来，随着多传感器深度融合及 5G 通导一体化网络的全面部署，高精度、低功耗、强鲁棒性将成为室内定位产业蓬勃发展的主旋律。

\begin{thebibliography}{9}
\bibitem{r1} 闫大禹, 宋伟, 王旭丹, 等. 国内室内定位技术发展现状综述[J]. 导航定位学报, 2019, 7(4): 5--12.
\bibitem{r2} 代阳坪, 刘天宇. 室内定位技术综述及发展前景探究[J]. 数字通信世界, 2024(08): 89--92.
\bibitem{r3} Jaiteg Singh, Noopur Tyagi, Saravjeet Singh, et al. A Systematic Review of Contemporary Indoor Positioning Systems: Taxonomy, Techniques, and Algorithms[J]. IEEE Internet of Things Journal, 2024, 11(21): 34717--34745.
\bibitem{r4} 王慧强, 高凯旋, 吕宏武. 高精度室内定位研究评述及未来演进展望[J]. 通信学报, 2021, 42(7): 165--179.
\bibitem{r5} 陈锐志, 郭光毅, 陈亮, 等. 室内高精度定位技术研究应用现状与发展趋势[J]. 武汉大学学报(信息科学版), 2023, 48(10): 1591--1600.
\bibitem{r6} 柳鑫. 基于 WiFi 指纹的室内定位算法研究与实现[D]. 南京信息工程大学, 2024.
\bibitem{r7} Yang Zheng, Min Sheng, Junyu Liu, et al. Exploiting AoA Estimation Accuracy for Indoor Localization: A Weighted AoA-Based Approach[J]. IEEE Wireless Communications Letters, 2019, 8(1): 65--68.
\bibitem{r8} Shuai He, Hang Long, Wei Zhang. Multi-antenna Array-based AoA Estimation Using Bluetooth Low Energy for Indoor Positioning[C]//2021 7th International Conference on Computer and Communications (ICCC). IEEE, 2021: 1782--1786.
\bibitem{r9} Mingli Zhu, Quan Li. 2D MUSIC algorithm based on fast dimensionality reduction[C]//2025 5th International Conference on High Performance Computing, Big Data and Communication Engineering (ICHBC). IEEE, 2025: 121--125.
\end{thebibliography}
\end{document}
